\documentclass[10pt,a4paper]{amsart}
\usepackage[margin=19mm]{geometry}
\usepackage{amsmath,amssymb,mathtools}
\usepackage[scr=rsfs,cal=euler]{mathalfa}
\usepackage{xfrac}
\usepackage[unicode,hidelinks,bookmarks=false]{hyperref}
\newcommand{\ie}{i.e.\ }
\newcommand{\cf}{cf.\ }
\newcommand{\resp}{respectively\ }
\newcommand{\Subsection}{\subsection}
\providecommand{\nobreakdash}{\nobreak\hbox{-}}
\setlength{\emergencystretch}{2em}
\multlinegap0pt
\usepackage{mathtools}
\usepackage{amssymb}
\usepackage{xfrac}
\usepackage{float}
\usepackage{dsfont}
\usepackage{xcolor}
\newtheorem{lemma}{Lemma}
\title{No-go theorem for probabilistic one-way secret-key distillation with zero error}
\author{Vishal Singh, Mark M. Wilde}
\date{Source: arXiv:2404.01392v2 (CC BY 4.0)}
\begin{document}
\maketitle

\noindent\emph{Source and licence.} No-go theorem for probabilistic one-way secret-key distillation with zero error. Version arXiv:2404.01392v2 (CC BY 4.0), distributed by arXiv under the Creative Commons Attribution 4.0 International licence, http://creativecommons.org/licenses/by/4.0/, as recorded in the arXiv licence metadata for this exact version and shown on the abstract page https://arxiv.org/abs/2404.01392v2. Full licence notice: \texttt{licenses/arxiv-2404.01392-CC-BY-4.0.txt}.\par\smallskip

\noindent\textit{Editorial scope.} Counted unit: the article's lemma lem:dephas\_ext\_form (source0.tex lines 465--476). The article's complete original proof of it is delivered with it (source0.tex lines 478--520, one \texttt{proof} environment of 43 lines). No mathematics is written, repaired or re-derived: the statement and the proof are the article's own lines, in the article's own notation, and no retained line is substituted or re-flowed.  Retained uncounted as the source-local context the proof needs: lines 438--440.  Nothing was omitted from any retained span, so the package carries no omission record.  No retained line contains a citation, so the wrapper carries no bibliography and no bibliography entry.  No retained line contains a figure environment, an image-inclusion command or drawing code, so no image is delivered and the package declares no asset dependency.  Editorial formatting only, in addition: an AMS \texttt{amsart} preamble that restates the article's own declarations and macros used by the retained lines, namely the environments \texttt{lemma} on separate counters, together with the standard packages those lines need.  The retained environments print in this document as Lemma 1 in order of appearance, whereas the article prints the same items with the numbering of its own full document.  The complete native source of the article as distributed in the arXiv e-print for this exact version is bundled unchanged as \texttt{sources/157051/source0.tex}.
\begin{equation}\label{eq:dephas_channel_defn}
        \Delta_{A}\!\left(\cdot\right) \coloneqq \Pi_A\left(\cdot\right)\Pi_A + [e]_A\left(\cdot\right)[e]_A,
\end{equation}

\begin{lemma}\label{lem:dephas_ext_form}
    Let $\omega_{AE}$ be an arbitrary state in the set $\mathcal{F}\!\left(\eta^{p,k}_{AB}\right)$, where $\eta^{p,k}$ is a doubly erased state of the following form:
    \begin{equation}
        \eta^{p,k}_{AB} = p~\gamma^k_{AB} + (1-p)[e]_A\otimes[e]_B,
    \end{equation}
    with $\gamma^k_{AB}$ a bipartite private state holding $\log_2 k$ secret key bits. Here systems $E$ and $B$ are isomorphic to each other.
    The action of the partially dephasing channel $\Delta_A$, defined in \eqref{eq:dephas_channel_defn}, on the state $\omega_{AE}$ results in the following state:
    \begin{equation}
    \Delta_A\!\left(\omega_{AE}\right) = p~\sigma_{AE} + (1-p)[e]_A\otimes\tau_E,
    \end{equation}
    where $\sigma_{AE} \in \mathcal{F}\!\left(\gamma^k_{AB}\right)$ and $\tau_E \in \mathcal{S}\!\left(E\right)$.
\end{lemma}

\begin{proof}
    For every quantum state $\omega_{AE} \in \mathcal{F}\!\left(\eta^{p,k}_{AB}\right)$, there exists a state $\omega_{ABE}$ with the following marginals;
    \begin{align}
        \operatorname{Tr}_E\!\left[\omega_{ABE}\right] &= \eta^{p,k}_{AB},\\
        \operatorname{Tr}_B\!\left[\omega_{ABE}\right] &= \omega_{AE}.
    \end{align}

    Recall the definition of the partially dephasing channel $\Delta_A$ from \eqref{eq:dephas_channel_defn}. 
    The state arising from the action of  $\Delta_A$ on the state $\omega_{ABE}$ can be decomposed as follows:
    \begin{equation}\label{eq:dephas_decomp}
        \Delta_{A}\!\left(\omega_{ABE}\right) = X_{ABE} + [e]_{A}\otimes Y_{BE},
    \end{equation}
    where $[e]_AX_{ABE} = X_{ABE}[e]_A = 0$. Since $\Delta_{A}\!\left(\omega_{ABE}\right)$ is a quantum state, $\Delta\!\left(\omega_{ABE}\right) \ge 0$, which implies the following operator inequalities:
    \begin{align}
        X_{ABE} &\ge 0,\label{eq:X_pos}\\
        [e]_A\otimes Y_{BE} &\ge 0,\\
        \Rightarrow Y_{BE} &\ge 0\label{eq:Y_pos},
    \end{align}
    where we have used the fact that $X_{ABE}$ and $[e]_A\otimes Y_{BE}$ are orthogonal to each other. 
    
    Consider the marginal of $\Delta_{A }\!\left(\omega_{ABE}\right)$ on systems $AB$, which is equal to the following:
    \begin{align}
        \operatorname{Tr}_{E}\!\left[\Delta_{A}\!\left(\omega_{ABE}\right)\right] &=\operatorname{Tr}_{E}\!\left[X_{ABE}\right] + \operatorname{Tr}_{E}\!\left[Y_{BE}\right]\otimes[e]_A.\label{eq:marg_A_decomp}
    \end{align}
    We can also evaluate the marginal by first tracing out the system~$E$ and then applying the partially dephasing channel; that is,
    \begin{align}
        \operatorname{Tr}_{E}\!\left[\Delta_{A}\!\left(\omega_{ABE}\right)\right] &= \Delta_A\left(\operatorname{Tr}_{E}\!\left[\omega_{ABE}\right]\right)\\
        &= p~\gamma^k_{AB} + (1-p)[e]_A\otimes[e]_B.\label{eq:marg_A_orig}
    \end{align}
    Comparing \eqref{eq:marg_A_decomp} and \eqref{eq:marg_A_orig}, and using the fact that $[e]_AX_{ABE} = 0$, we arrive at the following equalities:
    \begin{align}
        \operatorname{Tr}_{E}\!\left[X_{ABE}\right] &= p~\gamma^k_{AB},\label{eq:X_marg}\\
        \operatorname{Tr}_{E}\!\left[Y_{BE}\right] &= (1-p)[e]_B.\label{eq:Y_marg}
    \end{align}
    
    Recall from \eqref{eq:X_pos} and \eqref{eq:Y_pos} that $X_{ABE}$ and $Y_{BE}$ are positive semidefinite operators. Using \eqref{eq:X_marg}, we can further conclude that $\frac{X_{ABE}}{p}$ is a quantum state that extends $\gamma^k_{AB}$. Similarly, \eqref{eq:Y_marg} implies that $\frac{Y_{BE}}{1-p}$ is a quantum state that extends~$[e]_B$. 
    
    Any quantum state that extends $[e]_B$ is of the form $[e]_B\otimes\tau_E$, where $\tau_E$ is an arbitrary quantum state (and which can even be~$[e]_E$). Let us define the quantum state $\sigma_{AE} \coloneqq \frac{1}{p}\operatorname{Tr}_B\!\left[X_{ABE}\right]$. Since $\frac{1}{p}\operatorname{Tr}_E\!\left[X_{ABE}\right] = \gamma^k_{AB}$, $\sigma_{AE}$ is in the set $\mathcal{F}\!\left(\gamma^k_{AB}\right)$. The marginal of the state $\Delta_A\!\left(\omega_{ABE}\right)$ on systems $AE$ can now be written as follows:
    \begin{equation}
        \Delta_A\!\left(\omega_{AE}\right) = p~\sigma_{AE} + (1-p)[e]_A\otimes\tau_E,
    \end{equation}
    where $\tau_E \in \mathcal{S}\!\left(E\right)$ and $\sigma_{AE} \in \mathcal{F}\!\left(\gamma^k_{AB}\right)$.
\end{proof}

\end{document}
